% GENERATED FILE. Edit canonical lessons, then run npm run generate:books.
% canonical-source-hash: fnv1a64:b5d5b0c70654ec7d
% canonical-lessons: GE-C438-altstadt, GE-C438-kathedrale, GE-C438-glockenturm, GE-C438-stadtmauer, GE-C438-stadttor

\chapter{Die alte Stadt}
\label{ch:ge-die-alte-stadt}
\hlchaptermodality{438}

\begin{chapteropening}
\textbf{By the end of this chapter:} \emph{I can talk about the old part of a town: the old town, the cathedral, the bell tower, the town wall and the town gate.}

Complete the last lesson of chapter 438: I can talk about the old part of a town: the old town, the cathedral, the bell tower, the town wall and the town gate.
\end{chapteropening}

\section[die Altstadt]{die Altstadt — the old town}

\label{lesson:GE-C438-altstadt}



\begin{quote}
Before the new one: say the German for the information, then the German for to know one's way around.
\end{quote}

\subsection*{You'll want to know: die Altstadt}
\textbf{die Altstadt} — \textquotedblleft{}the old town\textquotedblright{}.

\begin{quote}
Die Altstadt hat viele kleine Gassen.
\end{quote}

\begin{quote}
\emph{The old town has lots of little lanes.}
\end{quote}

Plural: \emph{die Altstädte}, rare.

\begin{cousinweb}
\emph{Altstadt} is built from \emph{alt} 'old', inherited from Proto-Germanic \emph{aldaz} (EtymDB), and \emph{Stadt} 'town', from Old High German \emph{stat} 'place' (EtymDB).
\end{cousinweb}

\begin{grammarlens}[title={What you've built}]
One more word for telling, explaining and describing at B1.
\end{grammarlens}

\subsection*{Your turn}
\begin{itemize}
\raggedright
  \item \emph{Say it:} \emph{die Altstadt}
  \item \emph{Say it:} \emph{die Altstadt}, once more
  \item \emph{Say it:} a sentence of your own with \emph{die Altstadt}
  \item \emph{Recall:} say the German for the information, then the German for to know one's way around, then say \emph{die Altstadt} again
\end{itemize}

\begin{tcolorbox}[breakable,colback=teal!4,colframe=teal!35!black,title={Before you move on}]
What does die Altstadt mean? (\textquotedblleft{}The old town\textquotedblright{}.) Say it once more.
\end{tcolorbox}
\section[die Kathedrale]{die Kathedrale — the cathedral}

\label{lesson:GE-C438-kathedrale}



\begin{quote}
Before the new one: say the German for to know one's way around, then the German for the old town.
\end{quote}

\subsection*{You'll want to know: die Kathedrale}
\textbf{die Kathedrale} — \textquotedblleft{}the cathedral\textquotedblright{}.

\begin{quote}
Die Kathedrale ist über siebenhundert Jahre alt.
\end{quote}

\begin{quote}
\emph{The cathedral is over seven hundred years old.}
\end{quote}

Plural: \emph{die Kathedralen}. Many German towns say \emph{der Dom} or \emph{das Münster}.

\begin{cousinweb}
\emph{Kathedrale} is a learned borrowing from Medieval Latin \emph{(ecclēsia) cathedrālis} 'church with a bishop's chair', from Latin \emph{cathedra} 'chair', from Greek \emph{kathédra} (Lewis \& Short).
\end{cousinweb}

\begin{grammarlens}[title={What you've built}]
One more word for telling, explaining and describing at B1.
\end{grammarlens}

\subsection*{Your turn}
\begin{itemize}
\raggedright
  \item \emph{Say it:} \emph{die Kathedrale}
  \item \emph{Say it:} \emph{die Kathedrale}, once more
  \item \emph{Say it:} a sentence of your own with \emph{die Kathedrale}
  \item \emph{Recall:} say the German for to know one's way around, then the German for the old town, then say \emph{die Kathedrale} again
\end{itemize}

\begin{tcolorbox}[breakable,colback=teal!4,colframe=teal!35!black,title={Before you move on}]
What does die Kathedrale mean? (\textquotedblleft{}The cathedral\textquotedblright{}.) Say it once more.
\end{tcolorbox}
\section[der Glockenturm]{der Glockenturm — the bell tower}

\label{lesson:GE-C438-glockenturm}



\begin{quote}
Before the new one: say the German for the old town, then the German for the cathedral.
\end{quote}

\subsection*{You'll want to know: der Glockenturm}
\textbf{der Glockenturm} — \textquotedblleft{}the bell tower\textquotedblright{}.

\begin{quote}
Vom Glockenturm sieht man die ganze Stadt.
\end{quote}

\begin{quote}
\emph{From the bell tower the whole town can be seen.}
\end{quote}

Plural: \emph{die Glockentürme}.

\begin{cousinweb}
\emph{Glockenturm} is built from \emph{Glocke} 'bell', which goes back to Old High German \emph{glocca} (EtymDB) and is usually traced to a Celtic word like Old Irish \emph{clocc}, and \emph{Turm} 'tower', from Latin \emph{turris}.
\end{cousinweb}

\begin{grammarlens}[title={What you've built}]
One more word for telling, explaining and describing at B1.
\end{grammarlens}

\subsection*{Your turn}
\begin{itemize}
\raggedright
  \item \emph{Say it:} \emph{der Glockenturm}
  \item \emph{Say it:} \emph{der Glockenturm}, once more
  \item \emph{Say it:} a sentence of your own with \emph{der Glockenturm}
  \item \emph{Recall:} say the German for the old town, then the German for the cathedral, then say \emph{der Glockenturm} again
\end{itemize}

\begin{tcolorbox}[breakable,colback=teal!4,colframe=teal!35!black,title={Before you move on}]
What does der Glockenturm mean? (\textquotedblleft{}The bell tower\textquotedblright{}.) Say it once more.
\end{tcolorbox}
\section[die Stadtmauer]{die Stadtmauer — the town wall, the city wall}

\label{lesson:GE-C438-stadtmauer}



\begin{quote}
Before the new one: say the German for the cathedral, then the German for the bell tower.
\end{quote}

\subsection*{You'll want to know: die Stadtmauer}
\textbf{die Stadtmauer} — \textquotedblleft{}the town wall, the city wall\textquotedblright{}.

\begin{quote}
Um die Altstadt steht noch eine alte Stadtmauer.
\end{quote}

\begin{quote}
\emph{There is still an old wall round the old town.}
\end{quote}

Plural: \emph{die Stadtmauern}. \emph{Die Mauer} is a strong outside wall; \emph{die Wand} is a wall inside a house.

\begin{cousinweb}
\emph{Stadtmauer} is built from \emph{Stadt} 'town', an inherited word, and \emph{Mauer} 'wall', an early borrowing from Latin \emph{mūrus} (EtymDB; Lewis \& Short).
\end{cousinweb}

\begin{grammarlens}[title={What you've built}]
One more word for telling, explaining and describing at B1.
\end{grammarlens}

\subsection*{Your turn}
\begin{itemize}
\raggedright
  \item \emph{Say it:} \emph{die Stadtmauer}
  \item \emph{Say it:} \emph{die Stadtmauer}, once more
  \item \emph{Say it:} a sentence of your own with \emph{die Stadtmauer}
  \item \emph{Recall:} say the German for the cathedral, then the German for the bell tower, then say \emph{die Stadtmauer} again
\end{itemize}

\begin{tcolorbox}[breakable,colback=teal!4,colframe=teal!35!black,title={Before you move on}]
What does die Stadtmauer mean? (\textquotedblleft{}The town wall, the city wall\textquotedblright{}.) Say it once more.
\end{tcolorbox}
\section[das Stadttor]{das Stadttor — the town gate, the city gate}

\label{lesson:GE-C438-stadttor}



\begin{quote}
Before the new one: say the German for the bell tower, then the German for the town wall.
\end{quote}

\subsection*{You'll want to know: das Stadttor}
\textbf{das Stadttor} — \textquotedblleft{}the town gate, the city gate\textquotedblright{}.

\begin{quote}
Durch das alte Stadttor fahren keine Autos.
\end{quote}

\begin{quote}
\emph{No cars drive through the old town gate.}
\end{quote}

Plural: \emph{die Stadttore}.

\begin{cousinweb}
\emph{Stadttor} is built from \emph{Stadt} 'town' and \emph{Tor} 'gate', both inherited: \emph{Tor} goes back through Middle High German to Old High German \emph{tor} (EtymDB).
\end{cousinweb}

\begin{grammarlens}[title={What you've built}]
One more word for telling, explaining and describing at B1.
\end{grammarlens}

\subsection*{Your turn}
\begin{itemize}
\raggedright
  \item \emph{Say it:} \emph{das Stadttor}
  \item \emph{Say it:} \emph{das Stadttor}, once more
  \item \emph{Say it:} a sentence of your own with \emph{das Stadttor}
  \item \emph{Recall:} say the German for the bell tower, then the German for the town wall, then say \emph{das Stadttor} again
\end{itemize}

\begin{tcolorbox}[breakable,colback=teal!4,colframe=teal!35!black,title={Before you move on}]
What does das Stadttor mean? (\textquotedblleft{}The town gate, the city gate\textquotedblright{}.) Say it once more.
\end{tcolorbox}
